\documentclass[../main.tex]{subfiles}
\begin{document}

\chapter{Introducción}
\label{C:introduccion}

\section{Contexto y motivación}

La digitalización acelerada de los últimos años ha transformado los datos
en el activo más valioso de cualquier organización. Empresas, hospitales,
administraciones públicas e instituciones educativas almacenan hoy en día
volúmenes de información sin precedentes. Esta concentración de valor en
entornos digitales ha atraído inevitablemente la atención de los
ciberdelincuentes, que han desarrollado herramientas cada vez más
sofisticadas para explotar esa dependencia.

Entre todas las amenazas del panorama actual de ciberseguridad, el
ransomware ocupa un lugar destacado por su capacidad de causar daños
inmediatos, tangibles y difícilmente reversibles. Según el informe
IBM X-Force Threat Intelligence Index~2024~\cite{ibm2024}, el tiempo
medio que transcurre entre la intrusión inicial y la detección supera
los 200~días en entornos sin instrumentación activa. Las consecuencias
económicas son devastadoras: el informe de Veeam Data Protection
Trends~2024~\cite{veeam2024} estima que el coste medio de un ataque de
ransomware para una organización mediana supera los cuatro millones de
dólares.

El mecanismo de funcionamiento del ransomware moderno sigue un patrón
conocido como la cadena de ataque o \textit{kill chain}. Tras obtener
acceso inicial, el atacante pasa un período de reconocimiento en el que
mapea la red, escala privilegios y, crucialmente, localiza y destruye las
copias de seguridad. Solo entonces despliega el cifrador. Esta secuencia
es especialmente relevante para comprender por qué las estrategias de
backup tradicionales han ido perdiendo efectividad: los atacantes obtienen
con frecuencia credenciales que les permiten acceder y destruir los
repositorios antes del cifrado.

Frente a este panorama han surgido dos líneas de respuesta complementarias
que este trabajo estudia de forma conjunta: el almacenamiento inmutable
mediante Object Lock en plataformas S3-compatibles, y la instrumentación
del entorno con un \gls{siem} para detectar el ataque en las fases más
tempranas posibles.

La motivación personal de este proyecto surge directamente de la experiencia
profesional como administrador de sistemas. En el día a día de la
administración de infraestructuras IT, la decisión sobre qué solución de
backup adoptar raramente se apoya en datos empíricos comparativos. Este
trabajo intenta aportar datos concretos, obtenidos en un entorno
reproducible y con metodología estadística rigurosa.

\subsection{Marco regulatorio: presión normativa sobre la resiliencia
de datos}

La motivación técnica descrita anteriormente convive con una presión
regulatoria creciente que refuerza la relevancia práctica de este trabajo.
La Directiva (UE) 2022/2555, conocida como NIS2, entró en aplicación en
octubre de 2024 y amplía sustancialmente el número de sectores obligados a
implantar medidas de gestión de riesgos de ciberseguridad, incluyendo de
forma explícita la gestión de copias de seguridad y los planes de
recuperación ante desastres entre las medidas mínimas exigidas a las
entidades esenciales e importantes. El incumplimiento contempla sanciones
de hasta el 2\,\% de la facturación anual global para las entidades
esenciales.

En el ámbito español, el Esquema Nacional de Seguridad (ENS), regulado por
el Real Decreto 311/2022, establece para las administraciones públicas
requisitos específicos sobre continuidad del servicio (medida \texttt{mp.cont})
que incluyen la realización de copias de seguridad periódicas y pruebas de
restauración documentadas. Ambos marcos comparten un principio de fondo
directamente relevante para este \pfg: no basta con \emph{tener} una copia
de seguridad, sino que debe demostrarse que dicha copia sigue estando
disponible e íntegra frente a un adversario activo, no solo frente a fallos
accidentales de hardware. Esta distinción --- disponibilidad frente a un
fallo pasivo frente a disponibilidad frente a un adversario que
intenta activamente destruir la copia --- es precisamente el eje
experimental de este trabajo y se retomará con evidencia empírica en el
capítulo~\ref{C:discusion}.

Aunque el laboratorio desarrollado no pretende certificar el cumplimiento
de ninguna normativa concreta, los resultados obtenidos son directamente
trasladables a la justificación técnica que exige, por ejemplo, la medida
\texttt{mp.cont.2} (continuidad del servicio) del ENS o el artículo~21 de
NIS2 relativo a la gestión de la continuidad de negocio.

\subsection{Modelo de amenaza}
\label{sec:modelo-amenaza}

Para interpretar correctamente el alcance de los resultados presentados en los
capítulos siguientes conviene explicitar, ya desde la introducción, el modelo de
amenaza asumido por el laboratorio. Se adopta un modelo de adversario con las
siguientes capacidades y limitaciones, coherente con el comportamiento
documentado de los grupos \gls{raas} descrito en el
capítulo~\ref{C:estado_arte}:

\begin{description}
  \item[Capacidades asumidas] El atacante dispone de credenciales \gls{ssh}
  válidas sobre el \emph{fileserver}, obtenidas por un vector no modelado en el
  laboratorio (véase «Vector de acceso inicial» más abajo). A partir de ese punto,
  el atacante actúa con los mismos privilegios que un administrador de sistemas
  legítimo sobre los recursos accesibles desde esa cuenta: puede leer, modificar
  y eliminar ficheros del \emph{fileserver}, ejecutar \emph{scripts} arbitrarios,
  y utilizar las credenciales almacenadas localmente (claves \gls{ssh}, tokens de
  \gls{api}) para pivotar hacia otros sistemas alcanzables, como se modela en el
  escenario~E3.

  \item[Capacidades explícitamente excluidas] El atacante no dispone de acceso
  físico a ningún equipo, no puede comprometer el \emph{firmware} ni el
  \emph{hipervisor} subyacente, y no dispone de capacidad para modificar el
  binario o la configuración de MinIO ni de Wazuh a nivel de sistema operativo del
  \emph{host}. Esta última exclusión es la que sostiene la garantía de
  inmutabilidad evaluada en el escenario~E2: si el modelo de amenaza incluyera
  compromiso del \emph{host} de MinIO, la protección de Object Lock podría
  eludirse, como se reconoce explícitamente en la discusión del
  capítulo~\ref{C:discusion}.

  \item[Objetivo del atacante] Maximizar el daño mediante cifrado irreversible de
  los datos y, cuando es posible, la destrucción de las copias de recuperación
  disponibles, siguiendo el patrón de \emph{doble extorsión} descrito en la
  sección~\ref{sec:destruccion-backups}. El laboratorio no modela la fase de
  exfiltración con fines de extorsión adicional más allá del escenario~E4, que se
  limita a evaluar la detectabilidad de la transferencia de datos.

  \item[Nivel de sofisticación] Se asume un adversario con conocimientos técnicos
  equivalentes a un afiliado \gls{raas} de nivel medio: capaz de utilizar
  herramientas estándar de línea de comandos (\texttt{ssh}, \texttt{aws cli},
  \texttt{scp}) pero sin desarrollo de \emph{malware} a medida ni técnicas
  avanzadas de evasión de \gls{edr} como las descritas en la
  sección~\ref{sec:evasion-edr} del capítulo~\ref{C:estado_arte}. Esta
  simplificación es coherente con el alcance declarado en la sección siguiente y
  se revisita como limitación en el capítulo~\ref{C:discusion}.
\end{description}

\section{Objetivos}

El \textbf{objetivo general} del trabajo es evaluar de forma experimental
el impacto de distintas estrategias de \textit{backup} sobre la capacidad
de recuperación ante un ataque de ransomware simulado, en un entorno
instrumentado con un \gls{siem}.

De este objetivo general se derivan los siguientes \textbf{objetivos
específicos}:

\begin{enumerate}[label=OE\arabic*., leftmargin=2.5cm]
  \item Diseñar e implementar un laboratorio de simulación de ransomware
  basado en contenedores Docker que sea reproducible, portable y
  completamente automatizado.

  \item Implementar y comparar tres estrategias de \textit{backup}: ausencia
  de copia, copia en repositorio Restic sobre \gls{sftp} convencional, y
  copia en almacenamiento de objetos con Object Lock COMPLIANCE.

  \item Definir cuatro escenarios de ataque que cubran las tácticas más
  relevantes del ransomware moderno según la matriz MITRE ATT\&CK.

  \item Instrumentar el entorno con Wazuh \gls{siem} y desarrollar reglas
  de detección específicas para cada escenario, basadas en \gls{fim} y
  eventos estructurados de aplicación.

  \item Ejecutar un diseño factorial de 120~ejecuciones experimentales y
  analizar los resultados con pruebas estadísticas no paramétricas:
  Mann-Whitney~U, delta de Cliff, \textit{bootstrap} de intervalos de
  confianza y test exacto de Fisher.

  \item Validar las reglas de detección con técnicas reales del
  \textit{framework} MITRE ATT\&CK mediante Atomic Red Team.

  \item Publicar el laboratorio como software libre con documentación
  suficiente para reproducir los experimentos.
\end{enumerate}

\section{Hipótesis de investigación}

El trabajo plantea tres hipótesis:

\begin{description}
  \item[H1] La estrategia con Object Lock COMPLIANCE reduce el \gls{mttr}
  en al menos un 30\,\% respecto a la estrategia Restic sobre \gls{sftp}
  convencional.

  \item[H2] La detección automatizada mediante \gls{siem} permite
  identificar el inicio de un ataque en un tiempo mediano ($\tdetect$)
  inferior a 60~segundos en todos los escenarios evaluados.

  \item[H3] Ante la destrucción activa del repositorio de \textit{backup},
  la estrategia con Object Lock mantiene una tasa de recuperación exitosa
  significativamente superior a la estrategia SFTP (test exacto de Fisher,
  $\alpha = 0.05$).
\end{description}

\section{Alcance y limitaciones iniciales}

Para situar correctamente las aportaciones de este trabajo es necesario
delimitar explícitamente su alcance desde el planteamiento inicial, con
independencia de las limitaciones metodológicas que se detallarán tras la
presentación de resultados (capítulo~\ref{C:discusion}).

\begin{description}
  \item[Escala del laboratorio] El entorno reproduce una infraestructura
  empresarial \emph{mínima} (un fileserver, una workstation, un servidor de
  backup y un SIEM) desplegada sobre un único equipo físico mediante
  contenedores Docker. No se pretende representar la complejidad de una red
  corporativa real con segmentación VLAN, múltiples dominios de Active
  Directory o miles de endpoints.

  \item[Familias de ransomware simuladas] El laboratorio emplea un cifrador
  sintético desarrollado \textit{ad hoc} que reproduce el comportamiento
  esencial de un cifrador de ransomware (recorrido recursivo del sistema de
  ficheros, cifrado con AES-256-GCM y renombrado con extensión distintiva).
  No se ejecuta \textit{malware} real de ninguna familia conocida (LockBit,
  BlackCat, Cl0p, etc.), tanto por razones de seguridad del entorno de
  experimentación como por la imposibilidad de obtener licencias de
  investigación para muestras activas en el marco temporal de un \pfg.

  \item[Vector de acceso inicial] El experimento asume que el atacante ya
  dispone de acceso inicial al fileserver (mediante credenciales SSH
  comprometidas). No se modelan las fases previas de la cadena de ataque
  --- \textit{phishing}, explotación de vulnerabilidades públicas o compra
  de accesos en mercados clandestinos --- que en un incidente real
  preceden a la fase de impacto aquí estudiada.

  \item[Alcance del análisis SIEM] La instrumentación se centra
  deliberadamente en dos fuentes de datos (\gls{fim} y eventos de
  aplicación estructurados). Quedan fuera del alcance otras fuentes
  relevantes en un SOC real, como el análisis de tráfico de red mediante
  Suricata o Zeek, el análisis de comportamiento de usuario
  (\gls{ueba}) o la telemetría de \textit{endpoint} mediante \gls{xdr}
  comercial.

  \item[Generalización de resultados] Los resultados cuantitativos (MTTR,
  $\tdetect$) son específicos del hardware, la configuración de red y los
  parámetros de Wazuh empleados en este laboratorio concreto. Su valor
  principal no es la cifra absoluta sino el patrón cualitativo de
  disponibilidad diferencial entre estrategias, que se argumenta
  generalizable en el capítulo~\ref{C:discusion}.
\end{description}

\section{Estructura del documento}

El capítulo~\ref{C:estado_arte} revisa el estado del arte en ransomware,
estrategias de backup y sistemas \gls{siem}. El
capítulo~\ref{C:arquitectura} describe la arquitectura del laboratorio y
las decisiones de diseño. El capítulo~\ref{C:disenyo} detalla el diseño
experimental y la metodología estadística. El
capítulo~\ref{C:resultados} presenta y analiza los resultados. El
capítulo~\ref{C:discusion} discute las implicaciones de los resultados y
sus limitaciones. El capítulo~\ref{C:conclusiones} concluye el trabajo y
propone líneas de trabajo futuro.

Los anexos incluyen las reglas Wazuh completas (Anexo~A), la validación
con Atomic Red Team (Anexo~B) y la estructura del repositorio con
instrucciones de reproducción (Anexo~C).

\section{Planificación y presupuesto}

\subsection{Metodología de trabajo}

El desarrollo del proyecto siguió un enfoque iterativo de tipo
incremental, dividido en cuatro hitos verificables en lugar de un modelo
en cascada puro. Esta decisión responde a la naturaleza experimental del
trabajo: la definición final de las reglas de detección Wazuh y de los
umbrales de frecuencia solo pudo cerrarse tras observar el comportamiento
real del entorno durante la fase de calibración, lo que habría sido
imposible de anticipar completamente en la fase de diseño.

Los cuatro hitos fueron: (1)~\emph{arquitectura mínima viable}, con los
cuatro contenedores desplegados y comunicándose correctamente; (2)~
\emph{instrumentación}, con Wazuh desplegado y las reglas \gls{fim} básicas
operativas; (3)~\emph{automatización experimental}, con el
\textit{script} \texttt{experiment\_runner.py} capaz de ejecutar un
escenario completo de principio a fin sin intervención manual; y
(4)~\emph{ejecución y análisis}, con las 120~repeticiones completadas y el
análisis estadístico finalizado. Cada hito se cerró con una prueba de
regresión manual que verificaba que los hitos anteriores seguían
funcionando, dado que varios cambios en la configuración de red de Docker
durante la fase~3 afectaron a la conectividad establecida en la fase~1.

Como gestión de riesgos, se identificaron dos riesgos principales al
inicio del proyecto: la posible incompatibilidad de herramientas de
auditoría del sistema con el kernel del equipo de desarrollo (riesgo que
se materializó y se documenta en el capítulo~\ref{C:estado_arte}), y la
posibilidad de que el tiempo de ejecución de las 120~repeticiones excediera
el margen disponible antes de la fecha de entrega. Este segundo riesgo se
mitigó anticipando la automatización completa del experimento en la
fase~3, de modo que la ejecución final pudiera relanzarse sin supervisión
si fuera necesario repetir alguna celda.

\subsection{Planificación temporal}

Este \pfg tiene asignados 18~créditos ECTS, lo que equivale a
450~horas de trabajo según el factor de 25~horas por crédito del
\gls{ects}.

\begin{table}[!h]
  \centering
  \caption[Planificación temporal del TFG por fases.]
          {Planificación temporal del TFG por fases.}
  \label{tab:planificacion}
  \begin{tabularx}{\textwidth}{X C{2.5cm} C{2cm}}
    \toprule
    \rowcolor{tabla-cabecera}
    \textcolor{white}{\textbf{Fase}} &
    \textcolor{white}{\textbf{Período}} &
    \textcolor{white}{\textbf{Horas}} \\
    \midrule
    Revisión bibliográfica y estado del arte & Feb--Mar 2026 & 60 \\
    Diseño de la arquitectura               & Mar 2026      & 40 \\
    Implementación del laboratorio           & Mar--Abr 2026 & 90 \\
    Desarrollo de escenarios y playbooks     & Abr--May 2026 & 50 \\
    Calibración y validación                 & May 2026      & 40 \\
    Ejecución del experimento                & Jul 2026      & 20 \\
    Análisis estadístico y figuras           & Jul 2026      & 30 \\
    Validación con Atomic Red Team           & Jul 2026      & 20 \\
    Redacción de la memoria                  & May--Jul 2026 & 80 \\
    Revisión y entrega                       & Jul 2026      & 20 \\
    \midrule
    \textbf{Total}                           & \textbf{22 semanas} & \textbf{450} \\
    \bottomrule
  \end{tabularx}
\end{table}

\subsection{Presupuesto}

El presupuesto contempla recursos humanos, hardware e infraestructura y
software. Todo el software utilizado es de código abierto y gratuito,
por lo que el coste de licencias es cero.

\begin{table}[!h]
  \centering
  \caption[Resumen del presupuesto total del proyecto.]
          {Resumen del presupuesto total del proyecto.}
  \label{tab:presupuesto}
  \begin{tabularx}{\textwidth}{X C{3cm}}
    \toprule
    \rowcolor{tabla-cabecera}
    \textcolor{white}{\textbf{Categoría}} &
    \textcolor{white}{\textbf{Coste (€)}} \\
    \midrule
    Recursos humanos (450\,h a 25\,€/h + 10\,h tutor) & 11.850,00 \\
    Hardware e infraestructura (amortización)           &    401,00 \\
    Software y licencias                                &      0,00 \\
    Costes indirectos (15\,\%)                          &  1.837,65 \\
    \midrule
    \textbf{Total sin IVA}  & \textbf{14.088,65} \\
    IVA (21\,\%)            &          2.958,62  \\
    \midrule
    \textbf{Total con IVA}  & \textbf{17.047,27} \\
    \bottomrule
  \end{tabularx}
\end{table}

El presupuesto total asciende a \textbf{17.047,27~€ con IVA}, cifra que
refleja el coste real de contratación de un perfil técnico junior en
ciberseguridad para el desarrollo completo del proyecto.

\end{document}
